\newpage
\section{Diskussion}
\label{sec:Diskussion}

Im Folgenden werden die Messergebnisse anhand der theoretischen Erwartungen
eingeordnet und mögliche Abweichungen diskutiert.

\subsection{Messung der Wellenlänge}

Die getrennten Fits ergeben $\lambda_1=\vWaveOne$ für das Gitter mit
\qty[per-mode=reciprocal]{80}{\per\milli\meter} und $\lambda_2=\vWaveTwo$
für das Gitter mit \qty[per-mode=reciprocal]{100}{\per\milli\meter}.
Beide Werte liegen oberhalb der Referenzwellenlänge von
\qty{632.8}{\nano\meter} \cite[2]{eLas_HeNe}. Die relativen Abweichungen
betragen \vWaveOneDeviation{} beziehungsweise \vWaveTwoDeviation{}.
Beim ersten Gitter liegt die Referenz innerhalb des einfachen
Unsicherheitsintervalls. Die Abweichung beträgt
\vWaveOneReferenceSigma{} Standardunsicherheiten. Beim zweiten Gitter
beträgt sie \vWaveTwoReferenceSigma{} Standardunsicherheiten, sodass
hier eine größere Abweichung von der Referenz verbleibt.
Die reduzierten $\chi^2$-Werte von \vWaveOneChi{} und \vWaveTwoChi{}
liegen unter eins. Die Streuung der Messpunkte um die jeweilige Fitkurve
ist daher mit den angesetzten Unsicherheiten vereinbar. Wegen der nur drei
Freiheitsgrade je Fit ist daraus allein keine sichere Aussage über eine
Überschätzung der Unsicherheiten möglich.

Die Positionsunsicherheiten wurden mit \qty{0.5}{\milli\meter} angesetzt.
Zusätzlich kann die endliche Breite der Beugungsmaxima deren Lokalisierung
erschweren. Auch eine Abweichung vom angenommenen senkrechten Einfall
oder ein systematischer Fehler des Schirmabstands kann den ermittelten
Wellenlängenwert verschieben, ohne eine entsprechend große Streuung um
die Fitkurve zu erzeugen. Insbesondere beim zweiten Gitter sind solche
Beiträge als mögliche Erklärung zu berücksichtigen. Ihr Einfluss lässt
sich ohne zusätzliche Angaben nicht quantitativ beurteilen.

\subsection{Bestimmung der Polarisation}

Die Polarisationsmessung stimmt qualitativ mit dem erwarteten Verlauf nach
dem Malus-Gesetz überein. Der Fit beschreibt die gemessene
$\cos^2$-Abhängigkeit mit einem Maximum bei $\varphi_0=\vPolTheta$ und
einer nahezu vollständigen Auslöschung in der dazu orthogonalen Richtung.
Dies ist mit überwiegend linear polarisiertem Laserlicht aufgrund des Brewsterfensters
im Resonator vereinbar.

Auch quantitativ werden die Messwerte gut durch die Fitfunktion beschrieben.
Das reduzierte $\chi^2_\nu=\vPolChi$ liegt jedoch unter eins und legt nahe,
dass die Messunsicherheiten zu groß abgeschätzt worden sein könnten.
Eine mögliche Ursache ist eine großzügige Abschätzung der Leistungsunsicherheit.
Eine eindeutige Zuordnung zur Leistungsmessung ist allerdings nicht möglich,
da auch die Winkelunsicherheit in den Fit eingeht.

\subsection{Vermessung der TEM-Moden}

Sowohl die fotografischen Aufnahmen als auch die mit der Photodiode
aufgenommenen Leistungsprofile stimmen qualitativ mit den erwarteten
Verläufen überein. Die TEM$_{00}$-Mode zeigt ein gaußförmiges Profil,
während bei der TEM$_{10}$-Mode zwei Maxima mit einem Minimum dazwischen
zu erkennen sind.

Für die Grundmode liegt das reduzierte $\chi^2_\nu=\vZeroChi$ nahe
bei eins. Das gaußförmige Modell beschreibt die Messwerte damit unter
Berücksichtigung der angesetzten Positions- und Leistungsunsicherheiten
quantitativ zufriedenstellend. Ein Hinweis auf deutlich überschätzte
Messunsicherheiten ergibt sich für diese Messreihe nicht.

Das gemessene TEM$_{10}$-Profil weist eine leichte Asymmetrie auf.
Eine mögliche Ursache ist die erneute Ausrichtung der Schlitz-Photodioden-Anordnung
vor dieser Messreihe. Hinter der Streulinse ist der Strahl aufgeweitet und
divergent, sodass das Licht an verschiedenen Positionen unter unterschiedlichen
Winkeln auf die Anordnung trifft. Bei einer Verkippung könnten die beiden
Seiten des Profils daher unterschiedlich effizient erfasst worden sein.
Dies würde zu einer systematischen Verzerrung des gemessenen Profils führen.

Das reduzierte $\chi^2_\nu=\vFirstChi$ liegt dagegen deutlich über eins.
Die Abweichungen vom symmetrischen TEM$_{10}$-Modell sind größer, als die
angesetzten Messunsicherheiten erwarten lassen. Die beschriebene
Asymmetrie beziehungsweise weitere systematische Messbeiträge können
dazu beitragen. Auch unterschätzte Unsicherheiten sind möglich.
Der negative Fitwert $B_{10}=\vFirstB$ ist als empirischer Offset zu
verstehen und nicht als physikalisch negative Untergrundleistung.
Zusammen mit der erhöhten Fitabweichung begrenzt dies die quantitative
Aussagekraft des idealisierten Modells. Die angegebenen
Parameterunsicherheiten berücksichtigen diese zusätzlichen systematischen
Beiträge nicht.

\subsection{Stabilität bei verschiedenen Resonatorlängen}

Für die konkav-konkave Konfiguration war die theoretische obere
Stabilitätsgrenze von $L=\qty{280}{\centi\meter}$ aufgrund der begrenzten
Länge des Messaufbaus nicht erreichbar. Daher konnte lediglich qualitativ
bestätigt werden, dass der Laser innerhalb des untersuchten Längenbereichs
bis $L=\vStableLast$ stabil betrieben werden kann. Die dabei beobachteten
Schwankungen der Ausgangsleistung lassen sich durch die notwendige erneute
Justage nach jeder Änderung der Resonatorlänge erklären. Da die maximale
Leistung nur bei optimaler Ausrichtung der Resonatorspiegel und des Laserrohrs
erreicht wird, können Unterschiede in der Justage zu deutlichen
Leistungsschwankungen führen.

Für die plan-konkave Konfiguration liegt die theoretische Stabilitätsgrenze
bei $L=\qty{140}{\centi\meter}$. Die gemessene Leistung nimmt bei Annäherung
an diese Grenze von \qty{3.4}{\milli\watt} bei $L=\qty{132.5}{\centi\meter}$
auf \qty{0.3}{\milli\watt} bei $L=\qty{137.5}{\centi\meter}$ stark ab.
Oberhalb der Grenze wird bei $L=\qty{140.5}{\centi\meter}$ und
$L=\qty{160.0}{\centi\meter}$ keine Laserleistung mehr angezeigt.
Dies bestätigt die erwartete Stabilitätsgrenze qualitativ.
Die übrigen Leistungsschwankungen innerhalb des stabilen Bereichs sind
ebenfalls mit den Unterschieden bei der erneuten Justage vereinbar.

\subsection{Multimodenbetrieb und Frequenzspektrum}

Die bei angenommener Raumtemperatur abgeschätzte Dopplerbreite von etwa
\vDoppler\,\unit{\giga\hertz} ist deutlich größer als die gemessenen
Modenabstände von \vLongSpacingMin{} bis \vLongSpacingMax{}.
Diese liegen damit etwa eine Größenordnung unter der Dopplerbreite,
sodass mehrere longitudinale Moden innerhalb der verbreiterten Emissionslinie
liegen können. Der Multimodenbetrieb und die Beobachtung mehrerer
Schwebungsfrequenzen sind daher plausibel.

Die gemessenen Modenabstände nehmen mit zunehmender Resonatorlänge ab
und stimmen qualitativ mit dem erwarteten Verlauf $c/(2L)$ überein.
Die Differenzen zwischen den gemessenen Mittelwerten und den theoretischen
Werten in \autoref{tab:v61-mode-spacings} sind jedoch überwiegend größer
als die angegebenen Unsicherheiten. Dabei wurden bislang nur die
Unsicherheiten der Resonatorlängen berücksichtigt, nicht aber diejenigen
der gemessenen Frequenzen. Die Unsicherheiten der Differenzen sind daher
unvollständig. Zusätzliche Beiträge durch die Frequenzauflösung und die
Bestimmung der Peakpositionen könnten die beobachteten Abweichungen erklären.
Ob diese Beiträge hierfür ausreichen, lässt sich ohne eine Abschätzung
ihrer Größe nicht quantitativ beurteilen.

\newpage
